\section{目标化合物及相关结构}
\label{sec:identity}

\subsection{问题与范围}

比较合成条件前，先确定研究对象。现有资料包括含Zn目标相的原始报道、无Zn四方Ba--Y硅酸盐的历史研究，以及只能提供结构参照的外围化合物。前者把目标相写作$\mathrm{Ba_5Y_{12}Zn[O(SiO_4)]_8}$，并以单晶X射线衍射确定结构\cite{ababaikeri2024ba5y12zn}；后两类虽在组成或晶体学描述上接近，却不能替代含Zn目标相。本节依次核对原始化学式和出处、历史命名与现代结构重审，并说明现有结构与相区证据可以支持的结论。具体合成变量放在第~\ref{sec:condition-matrix}节。

这样可以把“获得了什么相”和“采用了什么条件”分开。表中为每条记录保留原文化学式、本文的排版形式、结构表征结果和待解决问题。只有化学式与直接结构证据同时指向目标相的研究，才可作为目标相的直接结构依据。相关谱系资料用于识别误指认、非整比和竞争相；结构相关化合物只提示需要核查的局部结构单元，不改变目标相的定义。

\subsection{目标相出处}

在本文核查的文献中，$\mathrm{Ba_5Y_{12}Zn[O(SiO_4)]_8}$只有一篇直接报道\cite{ababaikeri2024ba5y12zn}。该文采用开放体系高温溶液法，并以单晶X射线衍射鉴定晶体\cite{ababaikeri2024ba5y12zn}。这是目前唯一直接确认目标相的研究，不能与相关化合物拼成所谓“共识配方”。它支持目标化学式、路线类型和单晶结构身份，但不能证明跨实验室复现或批量相纯窗口，也未完全厘清氧计量表达、Zn位点及其与无Zn四方谱系的关系。

表~\ref{tab:identity-ledger}分列原文写法和本文排版。后者只统一符号、下标和括号，不把变组成式改成定组成式，也不依据元素守恒自行修改已发表的化学式。

\begin{table*}[t]
  \centering
  \caption{目标化合物及相关Ba--Y硅酸盐的组成与结构依据。表中按文献对象与目标相的关系分类；无Zn或结构相关研究均不构成含Zn目标相的独立复现。}
  \label{tab:identity-ledger}
  \scriptsize
  \setlength{\tabcolsep}{1.8pt}
  \begin{tabular}{P{1.9cm} P{1.9cm} P{1.15cm} P{2.25cm} P{2.3cm} P{2.2cm} P{2.3cm}}
    \toprule
    原文化学式/名称 & 本文采用的化学式 & 文献 & 主要结构证据 & 与目标相的关系 & 可得出的结论 & 尚待研究的问题\\
    \midrule
    \shortstack[l]{$\mathrm{Ba_5Y_{12}Zn}$\\$\mathrm{[O(SiO_4)]_8}$} & \shortstack[l]{$\mathrm{Ba_5Y_{12}Zn}$\\$\mathrm{[O(SiO_4)]_8}$} & \cite{ababaikeri2024ba5y12zn} & 单晶X射线结构鉴定\cite{ababaikeri2024ba5y12zn} & 目标化合物直接文献 & 含Zn目标相及其合成路线类型 & 独立复现、批量相纯窗口及与无Zn谱系的对应\\
    $\mathrm{Ba_{5+x}Y_{13}Si_8O_{41}}$（暂定） & 保留变量$x$ & \cite{kolitsch2009crystal,wierzbickawieczorek2017high} & 伴生相记录与系统相生成筛查\cite{kolitsch2009crystal,wierzbickawieczorek2017high} & 同类Ba--Y四方硅酸盐 & 历史命名及组成筛查结果 & 精确计量、结构模型及与目标相的关系\\
    \shortstack[l]{$\mathrm{Ba_xY_{26}Si_{16}}$\\$\mathrm{O_{71+x}}$} & 保留非整比式 & \cite{yamane2024synthesis} & 单晶非整比调制结构重审\cite{yamane2024synthesis} & 同类Ba--Y四方硅酸盐 & 历史式的现代结构解释 & Zn是否可进入该谱系\\
    \shortstack[l]{$\mathrm{Ba_xY_{26}Si_{16}}$\\$\mathrm{O_{71+x}}$陶瓷} & 同上 & \cite{yamane2024microstructure} & 批量相组成与显微结构分析\cite{yamane2024microstructure} & 同类Ba--Y四方硅酸盐 & 非整比谱系的批量表征需求 & 次生相、外来SiO$_2$与本征相区的区分\\
    \shortstack[l]{$\mathrm{Ba_5Y_{13}}$\\$\mathrm{[SiO_4]_8O_{8.5}}$} & \shortstack[l]{$\mathrm{Ba_5Y_{13}}$\\$\mathrm{[SiO_4]_8O_{8.5}}$} & \cite{gulay2024navigation} & EDS、CRED、同步辐射PXRD与中子精修\cite{gulay2024navigation} & 同类Ba--Y四方硅酸盐 & 独立确认的四方Ba--Y结构与相区参考 & 与含Zn目标相的位点对应\\
    $\mathrm{BaY_{16}Si_4O_{33}}$ & $\mathrm{BaY_{16}Si_4O_{33}}$ & \cite{motozawa2022bay16si4o33} & 单晶结构及孤立正硅酸根单元\cite{motozawa2022bay16si4o33} & 结构相关化合物 & Ba--Y正硅酸盐的结构参照 & 不具备与含Zn目标相同相的依据\\
    \bottomrule
  \end{tabular}
\end{table*}

\subsection{四方谱系重审}

历史名称的延续不能代替结构的延续。早期助熔生长研究把$\mathrm{Ba_{5+x}Y_{13}Si_8O_{41}}$记为暂定伴生相\cite{kolitsch2009crystal}；后续系统助熔研究把它纳入多变量成相筛查，但对象仍是无Zn的Ba--Y四方硅酸盐\cite{wierzbickawieczorek2017high}。现代单晶工作把相关名义组成重审为$\mathrm{Ba_xY_{26}Si_{16}O_{71+x}}$非整比调制结构\cite{yamane2024synthesis}。相应陶瓷研究表明，批量样品还必须检查次生相和研磨介质带入的SiO$_2$\cite{yamane2024microstructure}。因此，历史名称需要回到结构模型和批量相组成重新核对；上述研究均未直接观察精确的含Zn目标式。

独立相区研究结合EDS、CRED、同步辐射PXRD和中子精修，确认了四方$\mathrm{Ba_5Y_{13}[SiO_4]_8O_{8.5}}$\cite{gulay2024navigation}\linebreak ，为无ZnBa--Y硅酸盐提供了不同于暂定命名的结构参照。$\mathrm{BaY_{16}Si_4O_{33}}$的单晶结构则展示了含孤立正硅酸根单元的Ba--Y参照结构\cite{motozawa2022bay16si4o33}。它可帮助检查多面体和硅酸根的组织方式，但不是含Zn目标相的直接证据。本节据此确定目标相的化学式，并把其余记录用于说明相关结构和竞争相；这些记录可支持的论断范围见第~\ref{sec:evidence-distance}节。
